\section{Independent normalization and analytic Picard proofs}\label{inv:wb:analytic-picard-proofs}
\subsection{The cusp structure sheaf}

\subsubsection{structure-sheaf cohomology of the cusp fibre}\label{inv:wbproof:lem:cusp-structure-sheaf-audit}
The following working-edition argument proves the corresponding assertion of \cref{inv:A1.0}.


\begin{proof}
Let \(D_1,D_2,D_3\simeq\mathbb P^1\) be the double curves and \(P,Q\) the
two triple points, with the branch order fixed in the proof of
\cref{mapsI:top:thm:p0-integral-specialization}.  The normalization complex is
\[
0\longrightarrow\mathcal O_W\longrightarrow
 \nu_*\mathcal O_{\mathrm{dP}_6}\xrightarrow{\delta_0}
 \bigoplus_{k=1}^3\mathcal O_{D_k}\xrightarrow{\delta_1}
 \mathbb C_P\oplus\mathbb C_Q\longrightarrow0,
\tag{WBA.9.10}\label{inv:wb:eq:9.10}
\]
where \(\delta_0\) is the difference of the two branch restrictions and
\[
 \delta_1(g_1,g_2,g_3)
 =(g_1(P)-g_2(P)+g_3(P),
   g_1(Q)-g_2(Q)+g_3(Q)).
\]
For completeness, exactness is stalkwise.  At a generic double point it is
the fibre-product sequence
\[
0\to\mathbb C\{x,y,z\}/(xy)
 \to\mathbb C\{y,z\}\oplus\mathbb C\{x,z\}
 \to\mathbb C\{z\}\to0.
\]
At a triple point it is the augmented branch complex for
\(\mathbb C\{x,y,z\}/(xyz)\): three branch rings, three pairwise
intersection rings, and the common residue field, with alternating
differences.  The augmented two-simplex complex is exact; successively
lifting compatible power series proves exactness of the analytic rings.
Away from the singular locus the first arrow is an isomorphism.

The same formulas give an exact complex with \(\mathbb C\) in place of
\(\mathcal O\).  Since \(\mathrm{dP}_6\) is rational and each \(D_k\) is a
rational curve, the positive structure-sheaf cohomology of every term on
the right of \eqref{inv:wb:eq:9.10} vanishes.  Hypercohomology therefore reduces to
\[
 \mathbb C\xrightarrow{0}\mathbb C^3
 \xrightarrow{d_1}\mathbb C^2,
 \qquad
 d_1(a_1,a_2,a_3)=(s,s),\quad s=a_1-a_2+a_3.
\tag{WBA.9.11}\label{inv:wb:eq:9.11}
\]
Its cohomology has dimensions \(1,2,1\), proving \cref{inv:A1.0}.

For the constant-sheaf complex, the row of degree zero is again \eqref{inv:wb:eq:9.11},
the row of degree one vanishes, and the remaining degree-two row begins
with \(H^2(\mathrm{dP}_6;\mathbb C)\).  The comparison induced by
\(\mathbb C_W\hookrightarrow\mathcal O_W\) is the identity on \eqref{inv:wb:eq:9.11} and
zero on the positive rows of the target.  It follows that the degree-one comparison of \cref{inv:A1.0} is an
isomorphism, while in degree two it maps the dual-complex summand
\(\ker\nu^*\) isomorphically onto \(H^2(W,\mathcal O_W)\).

Choose the basis \(\gamma_W,u_W\in H^1(W;\mathbb Z)\) whose intrinsic
specializations are \(\gamma,u\), respectively; it exists and is unique by
\cref{mapsI:top:thm:p0-integral-specialization}.  Their cup product is nonzero,
because its specialization is
\(\gamma\wedge u\ne0\in H^2(F;\mathbb Z)\).  Its pullback to
\(\mathrm{dP}_6\) is zero since
\(H^1(\mathrm{dP}_6;\mathbb C)=0\).  The degree-two comparison therefore
sends it to a nonzero element of \(H^2(W,\mathcal O_W)\).  The images of
\(\gamma_W,u_W\) form a basis by the degree-one comparison of \cref{inv:A1.0}, and both sides of the cup-product assertion of \cref{inv:A1.0} are
one-dimensional.  Hence the cup-product assertion of \cref{inv:A1.0} is an isomorphism.
\end{proof}

\subsection{Direct images, base change, and the Picard line}

\subsubsection{structure-sheaf direct images}\label{inv:wbproof:thm:structure-direct-images-audit}
The following working-edition argument proves the corresponding assertion of \cref{inv:A1.0}.


\begin{proof}
The proper map has connected fibres and \(B\) is normal, so Stein
factorization gives \(f_*\mathcal O_X=\mathcal O_B\).

The map is flat.  A local parameter on \(B\) pulls back to a nonzero
element of the regular local domain \(\mathcal O_{X,x}\); if it vanished
identically near one point, analytic continuation on the connected
manifold would make \(f\) constant.  Thus \(\mathcal O_{X,x}\) is
torsion-free, hence flat, over the discrete valuation ring
\(\mathcal O_{B,f(x)}\).

Over \(B^\circ\), the fibres are two-dimensional complex tori and their
structure-sheaf cohomology dimensions are \(1,2,1\).  At the cusp the same
dimensions hold by \cref{wb:lem:cusp-structure-sheaf-audit}.  Grauert's
theorem therefore gives local freeness of ranks \(1,2,1\) and base change
on a neighbourhood of \(B^\circ\cup\{p_0\}\); the precise analytic statement is
\cite[Chapter~III, \S4, Theorem~4.12]{BS76}, originally due to
Grauert \cite{Grauert1960}.

At \(p_j\), use the finite cover \(J'_j\to N_j\).  If
\(\widetilde f:J'_j\to\Delta_j\) is the smooth torus family and
\(\pi(s)=s^{m_j}\), exactness of finite-group invariants in characteristic
zero gives
\[
 R^qf_*\mathcal O_{N_j}
 =\bigl(\pi_*R^q\widetilde f_*\mathcal O_{J'_j}\bigr)^{G_j}.
\tag{WBA.9.14}\label{inv:wb:eq:9.14}
\]
After averaging, a free eigenframe \(v\) may be chosen for each character
\(\nu\).  The invariant sections are exactly
\(s^rv\,\mathcal O_{D_j}\), where \(0\le r<m_j\) is determined by
\(\nu\zeta_j^r=1\).  Hence \eqref{inv:wb:eq:9.14} is free of ranks \(1,2,1\).
Higher direct images vanish by the fibre-dimension bound, so these are all
the nonzero direct images.  Since all direct images are locally free, the analytic
cohomology-and-base-change theorem, applied by descending induction from
\(q=2\), proves the base-change assertion of \cref{inv:A1.0} at the finite fibres as well.  In the notation of
\cite[Chapter~III, \S3, Theorem~3.4, Corollaries~3.5, 3.7, and 3.10]
  {BS76}.  Corollary~3.10 says that, for each fixed \(q\),
  freeness of \(R^p f_*\mathcal O_X\) for every \(p\ge q\) is equivalent to
  cohomological flatness in all dimensions \(p\ge q\).  Applying this
  successively with \(q=2,1,0\) gives the displayed base-change maps in every
  relevant degree.  This is a local statement on the base and so applies
  equally to the nonreduced scheme-theoretic fibres.

Write \(G=R^1f_*\mathcal O_X\).  We construct its distinguished section
without using a cellular collapse.  On \(B^\circ\), the monodromy-invariant
integral class \(\gamma\) maps to a holomorphic section
\(\gamma^{0,1}\) of \(G\).  It is nowhere zero there: if its image in
\(H^{0,1}(F)\) vanished, then
\(\gamma\in H^{1,0}(F)\cap H^1(F;\mathbb R)=0\).  On a finite cover it is
a constant invariant class and remains nonzero at \(s_j=0\), so it
descends across \(p_j\).

At the cusp, let \(j:\Delta^*\hookrightarrow\Delta\).  The natural
specialization morphism
\[
 R^1f_*\mathbb C\longrightarrow
 j_*R^1f_*\mathbb C|_{\Delta^*}
\]
is an isomorphism: away from \(0\) this is tautological, while at \(0\) it
is exactly the isomorphism from \(H^1(W;\mathbb C)\) onto the
\(T_0\)-invariant subspace proved in
\cref{mapsI:top:thm:p0-integral-specialization}.  Hence the invariant classes
\(\gamma,u\) extend as sections of \(R^1f_*\mathbb C\) on the whole disc.
The natural morphism induced by \(\mathbb C\hookrightarrow\mathcal O\),
\[
 R^1f_*\mathbb C\otimes\mathcal O_\Delta
 \longrightarrow G|_\Delta,
\]
sends them at the central point to the two classes in
\(H^1(W,\mathcal O_W)\).  These form a basis by
\cref{wb:lem:cusp-structure-sheaf-audit}.  Thus
\(\gamma^{0,1}\) extends across the cusp and is nowhere zero everywhere.
In particular \(\mathcal O_B\gamma^{0,1}\subset G\) is a line subbundle.

Relative duality applies to the proper flat Gorenstein map \(f\): \(X\) is
smooth, the base is a smooth curve, and \(\omega_{X/B}\) is invertible.
Using \cref{wb:lem:hodge-line-audit}, it gives
\[
 R^2f_*\mathcal O_X\simeq\mathfrak K^\vee
 \simeq\mathcal O_B(-1).
\tag{WBA.9.15}\label{inv:wb:eq:9.15}
\]
The relative theorem is the one of Ramis--Ruget--Verdier
\cite[pp.~261--283]{RRV71}; the absolute singular
dualizing-complex comparison is independently recorded in
\cite[Theorems~1--2 and Proposition~2]{RR70}.

The relative cup product
\[
 \bigwedge^2G\longrightarrow R^2f_*\mathcal O_X
\tag{WBA.9.16}\label{inv:wb:eq:9.16}
\]
is an isomorphism.  It is so on every smooth torus fibre.  At the cusp it
is so by the base-change assertion of \cref{inv:A1.0} and \cref{wb:lem:cusp-structure-sheaf-audit}.  At \(p_j\),
choose on the cover an eigenframe \(v_1,v_2\), with
\(v_1=\gamma^{0,1}\) invariant.  Its invariant exponent is \(r_1=0\).
The action on \(v_1\wedge v_2\) is dual to the Hodge character in (9.3),
so its invariant generator is \(s_j^{a_j}(v_1\wedge v_2)\).
The invariant generator for the second line is consequently
\(s_j^{a_j}v_2\), because \(0<a_j<m_j\).  Wedge product sends
\(v_1\wedge(s_j^{a_j}v_2)\) to exactly the invariant generator of the
right side of \eqref{inv:wb:eq:9.16}, with no wrap modulo \(m_j\).

Consequently
\(\det G\simeq R^2f_*\mathcal O_X\simeq\mathcal O_B(-1)\).  The quotient
of \(G\) by its trivial subbundle is \(\mathcal O_B(-1)\), and the
extension splits because
\[
 \operatorname{Ext}^1(\mathcal O_B(-1),\mathcal O_B)
 =H^1(B,\mathcal O_B(1))=0.
\]
This proves \cref{inv:A1.0} and identifies the trivial summand.
\end{proof}

\subsubsection{the Picard parameter and fibrewise restriction}\label{inv:wbproof:cor:picard-line-audit}
The following working-edition argument proves the corresponding assertion of \cref{inv:C3.1}.


\begin{proof}
The Leray spectral sequence and \(H^1(B,\mathcal O_B)=0\) identify
\[
 H^1(X,\mathcal O_X)=H^0(B,R^1f_*\mathcal O_X)=\mathbb C
\]
by \cref{wb:thm:structure-direct-images-audit}.  Since the verified integral
topology gives \(H^1(X;\mathbb Z)=H^2(X;\mathbb Z)=0\), the analytic
exponential sequence
\[
0\to\mathbb Z\to\mathcal O_X
 \xrightarrow{u\mapsto e^{2\pi iu}}\mathcal O_X^*\to1
\]
identifies \(H^1(X,\mathcal O_X)\) with \(\operatorname{Pic}(X)\) and
shows that every line bundle is topologically trivial.  Under Leray,
restriction is evaluation of a global section of
\(R^1f_*\mathcal O_X\), followed by base change.  The only global section
is a scalar multiple of the nowhere-zero section \(\gamma^{0,1}\), so
evaluation is injective at every point and has the stated image on a
smooth fibre.
\end{proof}


